\documentclass{presentation}

\title{Διανύσματα}
\subtitle{Πρόσθεση - Αφαίρεση Διανυσμάτων}
\author[Λόλας]{Κωνσταντίνος Λόλας}
\date{}

\begin{document}

\begin{frame}
  \titlepage
\end{frame}

\begin{frame}{Καλά τα διανύσματα, αλλά τι τα κάνουμε?}
  \only<1>{
    Στους αριθμούς μπορούμε να κάνουμε πράξεις! Αντίστοιχα...
  }
  
  \only<2>{
    Θα κάνουμε πράξεις και με τα διανύσματα!
  }
\end{frame}

\begin{frame}{Πράξεις με Διανύσματα}
\begin{itemize}
  \item Πρόσθεση Διανυσμάτων
  \item Αφαίρεση Διανυσμάτων
  \item Πολλαπλασιασμός Διανυσμάτων
  \item Διαίρεση Διανυσμάτων NOPE
  \item Πολλαπλασιασμός Διανύσματος με Αριθμό
  

  και τέλος
\end{itemize}
\end{frame}

\begin{frame}{Πώς προσθέτουμε?}
  \begin{itemize}
  \item διαδοχικά
  \item κανόντας παραλληλογράμμου
  \end{itemize}
\end{frame}

\begin{frame}{Πώς προσθέτουμε? Άρα}
  Αν $ΟΑ=\vec{α}$ και $ΟΒ=\vec{β}$ τότε
  \begin{align*}
  \overrightarrow{ΑΟ}+\overrightarrow{ΟΒ}=\overrightarrow{ΑΒ}
  \end{align*}
\end{frame}

\begin{frame}{Αφαίρεση... σαν το Κιλκίς - Δεν υπάρχει!}

\end{frame}

\end{document}